\par\Needspace{6\baselineskip}\medskip
\begin{center}\hypertarget{native-173-article126}{}\pdfbookmark[2]{126.}{article-126}126.\end{center}
The preceding method leads directly to a \emph{third}. As before, let the distances of the celestial body from the Earth at the first and last observations be used to determine the corresponding longitudes in the orbit together with the radius vectors. With the aid of the orientation of the orbital plane supplied by this calculation, let the longitude in the orbit and the radius vector be found from the middle observation. Then let the remaining elements be computed from these three heliocentric positions by solving the problem treated in Articles 82 and 83. This operation is independent of the times of observation. In this way the three mean anomalies and the daily motion become known, from which we can calculate the intervals of time between the first and second observations and between the second and third. The differences of these intervals from the true ones are taken as $X$ and $Y$.

This method would be less suitable whenever the heliocentric motion covers only a small arc. In such a case this determination of the orbit depends on quantities of the third order, as we already pointed out in Article 82, and therefore does not admit of sufficient precision. The slightest variations in the values of $x,y$ could produce very large variations in the elements, and hence also in the values of $X,Y$, and the latter could not be assumed proportional to the former. Whenever, however, the three positions span a considerable heliocentric motion, the method will certainly work very well, provided it is unaffected by the exceptions \SourceLeaf{174}explained in the preceding articles, which clearly must be taken into account here as well.
